% !PRINTED 436
\weberpage{436}
Wenn wir nun die Bedingung aufsuchen, dass $\Theta^4 = 1$ wird,
so haben wir $\Theta^2 = \Theta^{-2}$ zu setzen. Bilden wir also $\Theta^2$ und $\Theta^{-2}$
nach (10), so kann man zuerst versuchen, die Uebereinstimmung
herzustellen, indem man eine der drei Zahlen $\alpha, \beta, \gamma$, etwa $\gamma = 0$,
setzt; dann muss aber nach (6) und (7) noch eine zweite, etwa
$\beta = 0$, und die dritte $\alpha = -1$ sein. Dann wird aber $\Theta^2$ nicht
mit $\Theta^{-2}$ übereinstimmend. Also sind alle drei Grössen $\alpha, \beta, \gamma$
von Null verschieden. Setzen wir nun das zweite und dritte
Glied der ersten Zeile in den aus (10) abgeleiteten Substitutionen
$\Theta^2$ und $\Theta^{-2}$ einander gleich, so ergiebt sich, wenn die Factoren
$\gamma$ und $\alpha$ abgeworfen werden:
\[
\alpha(\varepsilon - \varepsilon^{-2}) = \beta(\varepsilon^{-1} - 1), \quad \gamma(\varepsilon^4 - 1) = \beta(\varepsilon^3 - \varepsilon)
\]
oder
\[
\alpha(\varepsilon^2 - \varepsilon^{-2}) = \beta(\varepsilon^4 - \varepsilon^{-4}), \quad \gamma(\varepsilon^2 - \varepsilon^{-2}) = \beta(\varepsilon - \varepsilon^{-1}),
\]
und hieraus, wenn $h$ einen unbestimmten Factor bedeutet:
\weq{11}{\begin{aligned} \alpha &= h(\varepsilon^4 - \varepsilon^{-4}) \\ \beta &= h(\varepsilon^2 - \varepsilon^{-2}) \\ \gamma &= h(\varepsilon - \varepsilon^{-1}). \end{aligned}}

Der Factor $h$ wird nach (7) aus der Gleichung
\weq{12}{-1 = h(\varepsilon + \varepsilon^2 + \varepsilon^4 - \varepsilon^{-1} - \varepsilon^{-2} - \varepsilon^{-4})}
bestimmt, und es ergiebt sich nach Bd. I, \S.\,171
\weq{13}{h = \frac{i}{\sqrt{7}} = \frac{-1}{\sqrt{-7}}.}

Das Vorzeichen von $\sqrt{7}$ hängt von der Wahl von $\varepsilon$ ab und
ist positiv, wenn z.\,B.
\[
\varepsilon = e^{\frac{2\pi i}{7}}
\]
genommen wird. Dann erhält man für $\alpha, \beta, \gamma$
\[
\alpha = \frac{-2\sin\dfrac{8\pi}{7}}{\sqrt{7}}, \quad \beta = \frac{-2\sin\dfrac{4\pi}{7}}{\sqrt{7}}, \quad \gamma = \frac{-2\sin\dfrac{2\pi}{7}}{\sqrt{7}}.
\]

Aus (11) und (12) ergeben sich noch die Formeln
\weq{14}{\begin{aligned}
\alpha\varepsilon + \beta\varepsilon^4 + \gamma\varepsilon^2 &= 1 \\
\alpha\varepsilon^{-1} + \beta\varepsilon^{-4} + \gamma\varepsilon^{-2} &= 1, \\
\alpha + \beta\varepsilon^{-1} + \gamma\varepsilon^2 &= 0 \\
\alpha + \beta\varepsilon + \gamma\varepsilon^{-2} &= 0; \\
\alpha\beta\gamma &= \tfrac{1}{7},
\end{aligned}}
%CONT>>
